\section{Related Works}
\label{sec:related_works}

Financial fraud detection has historically been tackled using static, batch-trained machine learning models. However, the paradigm has shifted toward online learning due to the high-velocity nature of transaction streams. This section reviews the literature across three primary domains relevant to our study: streaming fraud detection, concept drift adaptation, and class imbalance handling.

\subsection{Fraud Detection in Data Streams}
The application of streaming algorithms in fraud detection has gained significant traction. Author B et al. (2024) \cite{ref_b} proposed a real-time fraud detection architecture utilizing Kafka streaming and an ensemble of machine learning models. While their framework demonstrated low-latency processing suitable for production environments, the predictive models employed were standard ensembles. Consequently, the models suffered from severe performance degradation when classifying the extreme minority class (fraud), as standard algorithms inherently prioritize global accuracy over minority recall. Similarly, other studies focusing on auto-encoders and Support Vector Machines \cite{ref_a} have shown promise but typically rely on periodic offline retraining, making them highly susceptible to obsolescence when novel fraud tactics emerge.

\subsection{Concept Drift Adaptation}
To address the dynamic evolution of fraud behaviors, researchers have heavily integrated concept drift detectors into streaming pipelines. Author C et al. (2025) \cite{ref_c} introduced an incremental gradient boosting tree mechanism paired with Adaptive Windowing (ADWIN) to monitor performance degradation and trigger localized model updates. ADWIN mathematically bounds the probability of a true distribution shift, allowing the model to forget obsolete data efficiently. Despite its excellent responsiveness to concept drift, the core boosting algorithm in \cite{ref_c} was not equipped with a cost-sensitive objective function, leading to unacceptably low fraud recall in highly skewed real-world distributions.

\subsection{Class Imbalance Handling in Online Learning}
Class imbalance remains the most insidious challenge in streaming fraud detection. Traditional data-level approaches like SMOTE are computationally prohibitive in high-speed single-pass streams. Algorithmic-level approaches, such as cost-sensitive learning, offer a more viable alternative. Author D et al. (2026) \cite{ref_d} proposed a two-stage cost-sensitive learning framework for data streams experiencing both drift and imbalance. However, their reliance on a two-stage preprocessing architecture incurs significant computational overhead and memory consumption, rendering it impractical for high-throughput payment systems where transaction authorization must occur in milliseconds.

\subsection{Research Gap and Our Contribution}
Table \ref{tab:research_gap} summarizes the limitations of the state-of-the-art literature and highlights the specific research gaps filled by our proposed method. As demonstrated, although recent studies attempt to tackle concept drift using adaptive mechanisms \cite{ref_c}, they predominantly overlook extreme class imbalance. Conversely, methods attempting to solve both issues \cite{ref_d} introduce architectural latency that violates strict streaming constraints. 

To bridge this crucial gap, our research proposes the Cost-Sensitive Adaptive Random Forest (CS-ARF). By embedding an \textit{Implicit Online Oversampling} strategy directly into the node-splitting criteria of the Hoeffding Trees, our model seamlessly resolves both concept drift and extreme class imbalance in a single-pass streaming computation, eliminating the need for expensive two-stage architectures.

% Tabel Research Gap
\begin{table}[htbp]
    \centering
    \caption{Comparison of State-of-the-Art Literature and Our Proposed Method}
    \label{tab:research_gap}
    \resizebox{\textwidth}{!}{%
    \begin{tabular}{|l|l|c|c|p{6cm}|}
        \hline
        \textbf{Reference / Year} & \textbf{Proposed Method} & \textbf{Drift Adaptation?} & \textbf{Imbalance Handling?} & \textbf{Limitations / Research Gap} \\ \hline
        
        Author A et al. (2025) \cite{ref_a} & Ensemble-Based Fraud Detection & \ding{55} No & \ding{51} Yes & Model is trained offline (batch processing), making it highly susceptible to obsolescence when new fraud patterns emerge. \\ \hline
        
        Author B et al. (2024) \cite{ref_b} & Real-Time Fraud using Kafka \& ML Ensemble & \ding{51} Yes & \ding{55} No & Focuses on infrastructure latency but employs standard models that suffer from severe degradation on the minority class. \\ \hline
        
        Author C et al. (2025) \cite{ref_c} & Incremental Gradient Boosting Trees & \ding{51} Yes & \ding{55} No & Excellent at detecting drift, but exhibits low fraud recall due to the lack of cost-sensitive objective functions. \\ \hline
        
        Author D et al. (2026) \cite{ref_d} & Two-stage cost-sensitive learning & \ding{51} Yes & \ding{51} Yes & Utilizes a two-stage preprocessing architecture that incurs significant computational overhead, rendering it impractical. \\ \hline
        
        \textbf{Proposed Method} & \textbf{Cost-Sensitive Adaptive Random Forest} & \textbf{\ding{51} Yes} & \textbf{\ding{51} Yes} & \textbf{Eliminates two-stage computational overhead by utilizing Implicit Online Oversampling directly within the algorithm in a single pass.} \\ \hline
    \end{tabular}%
    }
\end{table}
